% ===== 原文 p.30 =====
% 说明：本文件首行内容紧接 tr04b.tex 末尾的"该投影用来指出："，是 p.29→p.30 的跨页句，
%       两文件之间不得插空行（否则句子会被断成两段）。
% [原文 p.30]
在 $R$ 的一条测地边界上指定一点，等价于在该测地线的半领口邻域边界上指定一点。由于当核心测地线长度趋于零时，领口及其边界收敛到尖点区域及其边界，我们就得到了 $\what{\cT}_{g,n}$ 定义之扩充的一种描述，使其包含"若干长度 $L_j$ 为零"的曲面。

标准尖点区域由变量 $w=e^{2\pi iz}$ 单值化。尖点区域边界 $\Im z=1/2$ 上的一点，对应于 $|w|=e^{-\pi}$ 上的一点；而鉴于因子 $e^{-\pi}$，圆周上的一点对应于原点处的一个切向量。变量 $w$ 在乘以一个模为 1 的复数（unimodular number）的意义下唯一；圆周与原点处切向量的这种等同是典范的。对 $S^1$ 的无穷小生成元 $-L_jt_j$，穿越测地线时位移向左（与图 3 的正向扭转、向右位移相对照），参考点的移动使得曲面内部位于其左侧，原点处的切向量顺时针旋转，而对偶的余切向量逆时针旋转（这是复直线 $\C$-线的正方向）。把各种等价性复合起来，尖点区域边界上一点的 $S^1$ 主丛等价于该穿刺点余切线中的一个非零向量。

% [Tutor Note: $w=e^{2\pi iz}$ 把 $\Im z=1/2$ 映到 $|w|=e^{2\pi i(x+iy)}$ 的模 $e^{-2\pi y}=e^{-\pi}$（$y=1/2$），这就是原文"$|w|=e^{-\pi}$"与"因子 $e^{-\pi}$"的来历；而 $z\mapsto z+1$ 对应 $w\mapsto w$，故尖点区域边界的圆周 $\{|w|=e^{-\pi}\}$ 与"过原点的切向量"（相差正的实缩放）典范同构。由此 $\psi$ 线丛（穿刺点处余切线的线丛）与尖点边界上变点的 $S^1$ 主丛在拓扑上是同一个对象——这正是 p.30 末句"$\what{\beta}_j$ 与 $\psi_j$ 拓扑等价"的来源，也是为什么 $\psi$ 类可以用 Chern 类来计算。]

\begin{proof}[定理 10 的证明]
$\MCG(R)$ 与 $(S^1)^n$ 在 $\what{\cT}_{g,n}$ 上的作用彼此交换；考察商 $\what{\cT}_{g,n}/\MCG(R)\times(S^1)^n$。由 Duistermaat--Heckman 定理 [CdS01, 第 30 章, 定理 30.8]，对 $\what{L}$ 的小值（包括 0），约化水平集 $\mu^{-1}(\what{L})/(S^1)^n$ 彼此微分同胚。进而由 Duistermaat--Heckman，$L$ 水平集的约化辛形式 $\omega_{\what{\cT}_{g,n}}\big|_{\mu^{-1}(\what{L})/(S^1)^n}$ 上同调于"0 水平集的约化形式 $\omega_{\what{\cT}_{g,n}}\big|_{\mu^{-1}(0)/(S^1)^n}$ 与诸贡献项 $(L_j^2/2)c_1(\what{\beta}_j)$ 之和"，其中 $c_1(\what{\beta}_j)$ 是第 $j$ 个尖点区域边界上变动的点的 $S^1$ 主丛的第一 Chern 类。结合命题 11，就得到所需的上同调等价，
\[
\omega_{\cT_g(L)}\equiv\omega_{\cT_g(0)}+\sum_{j=1}^n\frac{L_j^2}{2}c_1(\what{\beta}_j).
\]
由领口与穿刺点处余切线的描述，圆周丛 $\what{\beta}_j$ 在拓扑上等价于 $\psi$-线丛 $\psi_j$（见第 1 讲），且第一 Chern 类相等。证明完成。
\end{proof}

% [Tutor Note: 这段证明把 p.25 定理 10 的两个来源接了起来：Duistermaat--Heckman 给出"约化辛形式随矩值线性变化、斜率是第一 Chern 类"，于是有 $\omega_{\cT_g(L)}\equiv\omega_{\cT_g(0)}+\sum_j\frac{L_j^2}{2}c_1(\what{\beta}_j)$；再由末句 $\what{\beta}_j$ 与 $\psi_j$ 的拓扑等价、$c_1(\psi_j)=c_1(\what{\beta}_j)$，把它代回定理 10 的 $\frac1{d!}\int_{\cT_{g,n}/\MCG}(\omega+\sum_j\frac{L_j^2}{2}\psi_j)^d$，二项式展开后 $L$ 的幂次与总次数 $6g-6+2n=2d$ 恰好匹配。注意矩映射取的是 $\mu=\half L^2$，这正是求和项里出现 $L_j^2/2$（而不是 $L_j$）的原因；若按 $L_j$ 记矩值，公式会差一个因子 2，可作为校对锚点。]

\para{第 5 讲：相交数的模式与 Witten--Kontsevich。}

我们从 Harris--Morrison [HM98, 第 71--75 页] 的讨论开始。对有限个非负整数构成的序列 $\{\alpha_j\}$，定义最高次的 $\psi$-相交数为
\[
\langle\tau_{\alpha_1}\tau_{\alpha_2}\cdots\tau_{\alpha_n}\rangle_g=\int_{\Mgnb}\psi_1^{\alpha_1}\psi_2^{\alpha_2}\cdots\psi_n^{\alpha_n}.
\]
对非平凡的配对，亏格 $g$、穿刺点个数 $n$ 与指数 $\alpha_j$ 满足关系 $3g-3+n=\sum_{j=1}^n\alpha_j$；否则该配对定义为零。更一般地，用指数表示变量 $\tau$ 的幂（即重复），定义

% [Tutor Note: 本句自 p.30 末"…otherwise the pairing is defined as zero. More"起、至 p.31 首"generally using exponents to denote powers (repetitions) of the variables τ, define"止，跨页句由本块（tr-04）完整译出，译文到此为止；紧随其后的公式 $\langle\tau_0^{d_0}\tau_1^{d_1}\cdots\tau_m^{d_m}\rangle_g=\int_{\Mgnb}\prod_{j=0}^{m}\prod_{k=1}^{d_j}\psi_{(j,k)}^{j}$ 由 tr-05 的 tr05a.tex 承接，本文件不写该公式。读时注意两种记号的差别：上一式按穿刺点编号 $1,\dots,n$ 写 $\tau_{\alpha_i}$，下一式按幂次分组写 $\tau_j^{d_j}$（即 $d_j$ 个穿刺点各带 $\psi$ 的 $j$ 次方），两者在形式上不同但配对完全相同。]
